\item Con métodos experimentales particulares, es posible producir y observar en una larga barra delgada tanto una onda transversal cuya rapidez depende principalmente de la tensión en la barra como una onda longitudinal cuya rapidez está determinada por el módulo de Young y la densidad del material de acuerdo a la expresión $v = \sqrt{Y/\rho}$. La onda transversal se puede modelar como una onda en una cuerda estirada. Una barra metálica particular tiene 150 cm de largo y un radio de 0.200 cm y una masa de 50.9 g. El módulo Young para el material es $6.80 \times 10^{10}\text{ N/m}^2$. ¿Cuál debe ser la tensión en la barra si la razón de la rapidez de las ondas longitudinales a la rapidez de las ondas transversales es 8.00?